% =============================================================
%  CineMatch — Final Report (IEEE Conference Format)
%  EGN6933 Applied Data Science — University of Florida
%  Student: Roop Yaswanth Nagabhairava (45439754)
% =============================================================
\documentclass[conference]{IEEEtran}
\IEEEoverridecommandlockouts

\usepackage{cite}
\usepackage{amsmath,amssymb,amsfonts}
\usepackage{graphicx}
\usepackage{textcomp}
\usepackage{xcolor}
\usepackage{booktabs}
\usepackage{url}
\usepackage{hyperref}
\usepackage{array}
\usepackage{multirow}
\usepackage{algorithm}
\usepackage{algorithmic}
\usepackage{tikz}
\usepackage{placeins}
\usetikzlibrary{shapes.geometric,arrows.meta,positioning,calc,fit,backgrounds,patterns}

\hypersetup{colorlinks=true,linkcolor=blue,citecolor=blue,urlcolor=blue}
\renewcommand{\arraystretch}{1.15}

% ─────────────────────────────────────────────────────────────
\begin{document}

\title{CineMatch: A Cross-Cultural Movie Recommender via\\
Late-Fusion of Graph Neural Networks,\\Dense Multilingual Retrieval,
and Bayesian Quality Reranking}

\author{
  \IEEEauthorblockN{Roop Yaswanth Nagabhairava}
  \IEEEauthorblockA{M.S.\ Applied Data Science, University of Florida\\
    Student ID:~45439754 \quad egn6933@ufl.edu}
  \and
  \IEEEauthorblockN{Advisor: Dr.\ Yu Wang\quad Instructor: Dr.\ Edwin Marte}
  \IEEEauthorblockA{EGN6933: Applied Data Science, University of Florida}
}
\maketitle

% ─────────────────────────────────────────────────────────────
\begin{abstract}
Current recommendation platforms create \emph{cultural filter bubbles}:
collaborative filtering (CF) models produce no cross-cultural signal when
users' viewing histories are disjoint, and genre-tag content systems cannot
capture the narrative similarities that unite films across language boundaries.
CineMatch addresses this through a five-stage inference pipeline:
(1)~a multilingual BGE-M3 dense retrieval index (0.6\,B params, 1024-dim)
over 1.37\,M movies---merged from TMDB and IMDb metadata, refreshed
monthly via the TMDB API---indexed in a FAISS \texttt{IndexIDMap2(FlatIP)};
(2)~an XSimGCL graph neural network trained on MovieLens-32M (20.2\,M
positive interactions); both signals blended by (3)~an MLP routing gate
conditioned on five user-activity features producing mixture weights
$(\alpha,\beta,\gamma)$; candidates then pass through (4)~a
language-specific Bayesian quality filter and a genre-dominance reranker;
and (5)~a two-stage Determinantal Point Process (DPP) that enforces
cross-cultural diversity.
A state-of-the-art Qwen3-Embedding (4\,B) encoder was evaluated in
development but its 14\,GB index exceeded deployment constraints; BGE-M3
(5.6\,GB) was selected for production.
The XSimGCL backbone achieved \textbf{NDCG@10\,=\,0.7299},
\textbf{Recall@10\,=\,0.7277}, and \textbf{MRR@10\,=\,0.8045} on
979 held-out users (paired $t$-test vs.\ LightGCN: $t\!=\!12.65$,
$p\!<\!0.001$).
The full system is deployed at \url{onlymovies.app}---a Next.js frontend
on Vercel with a Dockerised FastAPI backend on Hugging Face Spaces---and
a live user study showed increase in
cross-cultural recommendation exposure over CF-only.
\end{abstract}

\begin{IEEEkeywords}
recommendation systems, graph neural networks, multilingual embeddings,
contrastive learning, Bayesian quality scoring, cross-cultural discovery,
determinantal point process, late fusion
\end{IEEEkeywords}

% ─────────────────────────────────────────────────────────────
\section{Problem Statement}


Most people's movie discovery is bounded by the language of their existing
watch history.
Streaming platforms optimise for engagement within familiar cultural
territory, reinforcing the standard CF paradigm: when users share no
cross-cultural viewing history, the model learns no cross-cultural signal.
Content-based genre tags compound the problem---``Action/Thriller'' does
not encode narrative structure, so a system cannot infer that a viewer
who loved \emph{The Dark Knight} might equally enjoy \emph{Vikram Vedha},
or that a Hayao Miyazaki fan might resonate with the structural language of
Guillermo del Toro.
Beyond content, the problem is structural: IMDb vote counts for regional
films are typically 10--100$\times$ lower than equivalent English-language
titles, so any quality filter using a flat vote threshold ($\geq3{,}000$
votes) silently suppresses the majority of world cinema before retrieval
even begins.

The global streaming market is projected to exceed \$230\,B by 2028, with
the majority of new subscribers coming from non-English-speaking markets.
Yet recommendation infrastructure overwhelmingly defaults to English-first
quality signals and English-centric interaction data.
A system capable of cross-cultural movie discovery---without requiring users
to already know what they are looking for---has both commercial value
(content monetisation across language barriers) and cultural impact
(preserving discoverability of non-dominant cinema).
Fifteen weeks of development were required to integrate three independently
deep research areas---graph contrastive learning, dense multilingual
retrieval, and probabilistic diversity---into a single production pipeline;
each required domain-specific engineering decisions documented here.

\textbf{The refined scope of CineMatch involves:}
building a hybrid retrieval system that identifies semantically similar
films across language and cultural boundaries; personalises for users
at every activity level from cold-start to power user via adaptive MLP
gating; applies language-aware quality scoring that treats regional cinema
equitably; enforces cross-cultural diversity via principled probabilistic
reranking; and operates at production latency in a fully deployed web
application.

% ─────────────────────────────────────────────────────────────
\section{Related Work}

\textbf{Matrix Factorisation foundations.}
ALS and SVD decompose the user-item rating matrix into latent factors
modelling only \emph{first-order} connectivity~\cite{movielens}.
They assign zero vectors to unrated items, lack any mechanism for
high-order neighbourhood signals, and embed popularity bias directly
into factor norms, trapping users in genre filter bubbles.
CineMatch includes ALS as a historical baseline to quantify the gap
(ALS: NDCG@10\,=\,0.4821; XSimGCL: 0.7299---a 51\% improvement).

\textbf{Graph neural network CF: LightGCN.}
LightGCN~\cite{lightgcn} simplified GCN-based CF by eliminating
non-linear transformations and learning only weighted neighbour
aggregation over the bipartite user-item graph.
While achieving strong recall on standard benchmarks, LightGCN suffers
from representation collapse on long-tail items: films with few
interactions receive weak propagation and converge to near-zero embeddings.
For cross-cultural recommenders, where non-English films are structurally
under-represented in training data, this is a critical failure mode
that CineMatch's contrastive backbone directly addresses.

\textbf{Graph contrastive learning: SimGCL and NCL.}
SimGCL~\cite{xsimgcl} addressed long-tail collapse by replacing
graph-topology augmentation with uniform noise perturbations on embeddings,
training with an InfoNCE contrastive loss; NCL~\cite{ncl} augments
contrastive learning with prototype-based neighbourhood signals.
Both significantly improve long-tail recall; however, SimGCL's approach
still applies noise through the full graph at training time.
CineMatch reproduces both as direct baselines
(SimGCL: NDCG@10\,=\,0.6920; NCL: 0.6890 on ML-32M).

\textbf{XSimGCL: our CF selection.}
Yu~et~al.~\cite{xsimgcl} showed that applying uniform noise to a
\emph{single} intermediate GCN layer---without any message-passing in
the contrastive view---achieves superior NDCG at lower training cost
than SimGCL.
CineMatch adopts XSimGCL as its CF backbone, confirmed by our
ML-32M results (NDCG@10\,=\,0.7299 vs.\ SimGCL 0.6920).

\textbf{Multilingual dense retrieval: Qwen3-Embedding and BGE-M3.}
Qwen3-Embedding (4\,B params) leads current multilingual retrieval
benchmarks (MTEB).
We encoded the full 1.37\,M-movie catalog with Qwen3, producing a
14\,GB FAISS index that outperformed BGE-M3 on niche semantic queries;
however, it exceeded Hugging Face Spaces' memory tier.
BGE-M3~\cite{bgem3} (0.6\,B params, 1024-dim, 100+ languages) was
selected for production: a 5.6\,GB index with comparable multilingual
alignment.


\textbf{Sequential / session-based models: SASRec and BERT4Rec.}
SASRec~\cite{sasrec} applies causal (left-to-right) self-attention
over a user's interaction sequence to predict the immediately \emph{next}
item, while BERT4Rec~\cite{bert4rec} uses bidirectional masked-item
prediction as a pre-training objective, allowing every past interaction
to attend to every other.
Both models excel when the \emph{current session context}---what the
user is in the mood for right now---differs meaningfully from their
long-run preferences.
CineMatch's goal is \emph{general taste exploration} rather than
next-item prediction within a bounded session; users visit the platform
to discover films across languages and eras, not to queue the next
autoplayed episode.
SASRec and BERT4Rec are therefore evaluated as related work but rejected
as the primary CF backbone: their sequential inductive bias would
discard historical preferences older than the session window and would
not leverage the long-tail cross-cultural signals captured by
XSimGCL's graph-level propagation.

\textbf{LLM + RAG hybrid retrieval.}
RAG-based recommenders combine a dense retrieval index (the ``retrieval''
stage) with a large language model that re-reads retrieved context to
produce either ranked scores or natural-language explanations~\cite{rag}.
For cold-start scenarios the LLM can synthesise a recommendation
from minimal profile text, and the approach inherently produces
human-readable rationales (``because you liked \emph{Parasite}, here
is \emph{The Wailing}'').
CineMatch's BGE-M3 path is architecturally equivalent to the
retrieval stage of an LLM+RAG pipeline: a query embedding is matched
against a pre-built dense index.

\textbf{Feature interaction models: DeepFM.}
DeepFM~\cite{deepfm} fuses a factorisation machine (FM) component for
low-order feature interactions with a deep neural network for
high-order interactions, operating over a rich tabular feature set:
user age, explicit genre tags, genome tag scores, and interaction
counts.
FM-family models excel on feature-rich datasets or proprietary platforms with
demographic metadata.
CineMatch trains and evaluates on MovieLens-32M, which intentionally
ships \emph{only} user–item rating triples without genome scores,
user demographics, or structured tag vectors.
Applying DeepFM would require imputing the absent tabular metadata,
introducing fabricated features that would corrupt ground-truth
correlations and make offline metrics unreliable.

\textbf{DPP for recommendation diversity.}
Chen~et~al.~\cite{dpp} demonstrated fast greedy MAP inference for
Determinantal Point Processes in recommendation, providing \emph{global}
diversity guarantees by maximising the determinant (geometric volume)
of the selected item set in $O(K^2n)$ time.
Compared to Maximal Marginal Relevance (MMR), which is a greedy
sequential heuristic prone to local optima, DPP accounts for all
pairwise item interactions simultaneously.
CineMatch extends this with a two-stage pipeline (categorical DPP
then embedding DPP), a novel design not present in prior work.

% ─────────────────────────────────────────────────────────────
\section{Objectives}

The architecture underwent four significant pivots.

\textit{(1) ALS $\to$ LightGCN (Weeks 4--7).}
Initial CF used ALS on MovieLens-1M.
Profiling showed ALS assigned near-zero vectors to items with fewer than
15 interactions---covering $>$60\% of non-English films---making it
unsuitable for cross-cultural objectives.
GNNs propagate neighbourhood signals over multiple hops, amplifying
long-tail items.

\textit{(2) LightGCN $\to$ UltraGCN $\to$ XSimGCL (Weeks 7--9).}
After LightGCN's NDCG@10\,=\,0.6410 on ML-32M, UltraGCN was evaluated
(0.6987); its constraint loss skips message-passing efficiently but
loses multi-hop signals on dense ML-32M.
XSimGCL (0.7299) was selected.

\textit{(3) Hard genre filter $\to$ soft dominance reranking (Week 11).}
Live users reported genre collapse after $\approx$40 interactions.
Root cause: a hard genre filter combined with a flat 3,000-vote floor
exhausted the candidate pool.
Genre preference was redesigned as a soft dominance score; vote thresholds
were made language-specific.

\textit{(4) Qwen3-4B $\to$ BGE-M3 for production (Week 12).}
The 14\,GB Qwen3 index exceeded Hugging Face Spaces' memory tier.
BGE-M3 (5.6\,GB) was deployed; Qwen3 remains available for future
infrastructure upgrades.


\begin{table}[h]
\caption{Project Objectives}
\label{tab:objectives}
\centering\small
\begin{tabular}{p{5.4cm}c}
\toprule
\textbf{Objective} & \textbf{Status} \\
\midrule
MovieLens-32M training & \checkmark \\
TMDB+IMDb merged catalog (1.37M movies) & \checkmark \\
Hybrid recommender (semantic + CF + reranking) & \checkmark \\
Cross-lingual discovery via multilingual embeddings & \checkmark \\
Bayesian quality scoring for regional cinema & \checkmark \\
Two-stage DPP diversity reranking & \checkmark \\
Live co-occurrence graph from user interactions & \checkmark \\
Deployed production web application & \checkmark \\
Statistical validation over baselines & \checkmark \\
Qwen3-4B encoder (SOTA) & Dev only* \\
RAG ``why this movie'' explanation layer & Descoped** \\
\bottomrule
\multicolumn{2}{l}{\small *14\,GB index exceeds production memory.}\\
\multicolumn{2}{l}{\small ** Listed as future work.}
\end{tabular}
\end{table}


% ─────────────────────────────────────────────────────────────
\section{Final Deliverables}

\begin{enumerate}
  \item \textbf{XSimGCL model}: trained on ML-32M; exported 512-dim
    user/item embeddings (200,807 and 65,023 rows respectively).
  \item \textbf{FAISS indexes}: BGE-M3 1024-dim, 1,366,255 movies
    (\texttt{IndexIDMap2(IndexFlatIP)}, 5.6\,GB, deployed);
    Qwen3-4B 14\,GB index (dev).
  \item \textbf{Python inference backend} (\texttt{app.py},
    $\approx$5,800 lines): initially validated via Gradio
    prototype; refactored to a \textbf{Dockerised FastAPI container}
    on Hugging Face Spaces.
    Implements the full pipeline: BGE-M3 retrieval, XSimGCL CF
    dot-product scoring, MLP gating, co-occurrence lookup,
    Bayesian quality filtering, language/genre reranking, two-stage DPP.
  \item \textbf{Next.js frontend} (\texttt{cinematch-web/}): deployed
    at \url{onlymovies.app} via Vercel; mobile-optimised swipe-based
    interaction logging to MongoDB.
  \item \textbf{MongoDB Atlas}: user profiles, session state,
    interaction logs ($\in\{\text{like,dislike,skip}\}$),
    \texttt{item\_cooccurrence} collection updated nightly.
  \item \textbf{GitHub}: \url{https://github.com/roop-yaswanth/CineMatch}\\
    (\texttt{README.md}, \texttt{requirements.txt}, training notebooks,
    FAISS build manifests, architecture reference).
\end{enumerate}

% ─────────────────────────────────────────────────────────────
\section{System Architecture}

The CineMatch inference pipeline is structured as a five-layer adaptive retrieval-reranking system (Fig.~\ref{fig:arch}). The logic moves from initial multi-path retrieval to adaptive modality fusion, followed by a multi-stage sequential reranking layer that enforces quality, diversity, and cultural alignment.

\begin{figure}[!t]
\raggedright
\includegraphics[width=0.48\textwidth]{architecture.jpg}
\caption{CineMatch five-layer inference architecture. The end-to-end pipeline begins when a request arrives as either a user identifier or a natural-language query. Layer~2 contains three parallel retrieval paths gated by an MLP that reads five user-activity features and outputs mixture weights $(\alpha,\beta,\gamma)$, blending BGE-M3 semantic search ($\alpha$), XSimGCL collaborative filtering ($\beta$), and live co-occurrence ($\gamma$). Cold-start users force $(\alpha,\beta,\gamma)\!=\!(1,0,0)$, while warm users receive an adaptive blend. The fused candidate pool then passes through Layer~4, which applies sequential Bayesian quality filtering, language-aware reranking, genre-dominance reranking, and a two-stage DPP diversity pass before final output.}
\label{fig:arch}
\end{figure}



% ─────────────────────────────────────────────────────────────
\FloatBarrier
\section{Methodology}
\label{sec:methodology}

\subsection{Data Pipeline}

\textbf{TMDB Catalog.}
A Kaggle-mirrored TMDB dump (1,367,793 entries) was filtered to remove
short films (runtime $<60$\,min) and TV episodes.
Final embeddable catalog: 1,366,255 records.
TMDB's \texttt{vote\_count} is unreliable for non-English films.
\texttt{merge\_imdb\_ratings.py} joined IMDb \texttt{numVotes} and
\texttt{averageRating} on \texttt{imdb\_id}, covering $\approx$738K titles.
A \texttt{best\_rating}/\texttt{best\_votes} composite column chooses
the more credible source per film.
The catalog is refreshed monthly via the TMDB API.

\textbf{MovieDoc Construction.}
Each movie becomes a structured text document fed to BGE-M3.
The \texttt{movieDoc} schema consists of:
\begin{itemize}
  \item \textbf{Title}: Primary film name.
  \item \textbf{Release date}: Calendar premiere date.
  \item \textbf{Year}: Release year (numeric).
  \item \textbf{IMDb ID}: IMDB identifier.
  \item \textbf{Original language}: Source cultural code.
  \item \textbf{Genres}: Categorical thematic tags.
  \item \textbf{IMDb rating}: Crowd-sourced score.
  \item \textbf{IMDb votes}: User interaction volume.
  \item \textbf{Best rating}: Merged quality signal.
  \item \textbf{Best votes}: Merged density signal.
  \item \textbf{Vote average}: TMDB-specific scoring.
  \item \textbf{Vote count}: TMDB interaction count.
  \item \textbf{Popularity}: Relative traffic score.
  \item \textbf{Tagline}: Condensed marketing hook.
  \item \textbf{Plot}: Deep narrative summary.
\end{itemize}
The \texttt{Best\_rating}/\texttt{Best\_votes} composite selects the
more credible source per film (IMDb where available, TMDB otherwise).
Combining the original language code, multiple vote signals, and full
plot lets BGE-M3 capture both narrative content and cultural identity.

\textbf{MovieLens-32M.}
ML-32M~\cite{movielens} (32M raw ratings) was filtered at
$\geq3.5$ stars to extract positive interactions:
200,807 users $\times$ 65,023 items, 20,218,546 interactions
(sparsity $= 99.85\%$).

\subsection{Semantic Retrieval --- BGE-M3 and Encoder Selection}

Three FAISS indexes were built and evaluated:
(1)~\textbf{BGE-M3} (0.6\,B params, 1024-dim): 1.37M movies, 5.6\,GB,
deployed in production.
(2)~\textbf{IMDb-subset BGE-M3}: 738K high-quality titles only.
(3)~\textbf{Qwen3-Embedding} (4\,B params): SOTA multilingual benchmark
scores; 14\,GB FAISS index, exceeds Hugging Face Spaces' memory tier.
BGE-M3 over the full catalog was selected as the optimal
NDCG-vs-memory trade-off.

BGE-M3 embeds each \texttt{movieDoc} to a 1024-dim $L_2$-normalised
vector: $\mathbf{e}_i = \text{BGE-M3}(\text{movieDoc}_i)$.
At query time, the user history is encoded with an instruction prefix:
\emph{``Instruct: Given liked movies, recommend similar ones.
Query: Title1 (Genre), Title2 (Genre)\ldots''}
Retrieval uses FAISS \texttt{IndexIDMap2(IndexFlatIP)} via
$s_\text{sem}(i)=\mathbf{q}^\top\mathbf{e}_i$ (cosine on normalised vectors).
The IDMap2 wrapper stores TMDB integer IDs directly in the index,
eliminating a secondary lookup and removing ID-mapping risks.
FlatIP (exact search) was preferred over IVF/HNSW approximation
to avoid accuracy loss from quantisation.
Index construction took 420\,s on an NVIDIA A100.

\subsection{Collaborative Filtering --- XSimGCL}

\textbf{Architecture.}
XSimGCL~\cite{xsimgcl} runs $L$ LightGCN propagation layers then
creates contrastive views by injecting uniform noise into a single
intermediate layer:
\begin{equation}
  \tilde{\mathbf{e}}_u^{(\ell)} = \mathbf{e}_u^{(\ell)} + \boldsymbol\delta_u,
  \quad\boldsymbol\delta\sim\mathcal{U}[-\varepsilon,\varepsilon].
\end{equation}
The InfoNCE contrastive loss pulls perturbed views of the same node
together while pushing apart different nodes:
\begin{equation}
  \mathcal{L}_\text{cl} = -\!\sum_u \log
    \frac{e^{\mathrm{sim}(\mathbf{e}_u,\tilde{\mathbf{e}}_u)/\tau}}
         {\sum_{u'\neq u}e^{\mathrm{sim}(\mathbf{e}_u,\tilde{\mathbf{e}}_{u'})/\tau}},
\end{equation}
\begin{equation}
  \mathcal{L} = \mathcal{L}_\text{BPR} + \lambda_\text{cl}\mathcal{L}_\text{cl}
              + \lambda_\text{reg}\|\Theta\|^2.
\end{equation}
The five training steps are: (1)~bipartite graph ingestion;
(2)~LightGCN multi-hop message passing; (3)~uniform noise augmentation
of one layer; (4)~BPR + InfoNCE loss computation; (5)~weighted sum for
weight update.

\textbf{Hyperparameters.}
Embedding size 512; $L=3$; noise $\varepsilon=0.1$; CL rate
$\lambda_\text{cl}=0.5$; temperature $\tau=0.2$;
$\lambda_\text{reg}=10^{-4}$; Adam lr$=10^{-3}$; batch 262,144;
mixed-precision on.
Training ran 47 effective epochs (early stop on NDCG@20) in
6,892\,s ($\approx$1.9\,h) on an NVIDIA B200 (HiPerGator cluster).

\textbf{Inference.}
$s_\text{CF}(i) = \mathbf{u}^\top\mathbf{e}_i$ via batched dot product
over 65,023 item embeddings.

\subsection{Cold-Start CF Embedding}

For users with $d_u<5$ interactions, a proxy 512-dim embedding is
synthesised by recency-weighted mean-pooling of liked items' XSimGCL
embeddings:
\begin{equation}
  \mathbf{u}_\text{cold} =
    \frac{\sum_{i\in\mathcal{L}} w(t_i)\,\mathbf{e}_i}
         {\sum_{i\in\mathcal{L}} w(t_i)},\quad
  w(t)=e^{-\lambda_d(T-t)},
\end{equation}
where $\lambda_d=\ln2/90$ (90-day half-life).
This is blended with a demographic cluster centroid at weight
$\gamma_d=0.30$:
$\hat{\mathbf{u}} = 0.7\,\mathbf{u}_\text{cold} + 0.3\,\mathbf{c}_\text{demo}$.
Centroids are looked up by hierarchical fallback:
(age, gender, region) $\to$ (age, gender) $\to$ (age) $\to$ global.

\subsection{MLP Gating and Late Fusion}

Five user-activity features drive the gate:
$\mathbf{f}_u = [\log d_u,\;\sigma^2_r,\;\rho,\;H_\ell,\;c_\ell]$
(degree, rating variance, recent-interaction fraction, language
entropy, cross-lingual ratio).
A two-layer MLP maps these to mixture weights:
\begin{equation}
  (\alpha,\beta,\gamma) =
  \mathrm{Softmax}(W_2\,\mathrm{ReLU}(W_1\mathbf{f}_u+b_1)+b_2).
  \label{eq:gate}
\end{equation}
Weights were hand-initialised from ablation experiments; end-to-end
gradient training is feasible once interaction logs accumulate.
Cold-start users are hard-routed: $(\alpha,\beta,\gamma)=(1,0,0)$.

\textbf{Modality Collapse Prevention.}
The two towers never share vector space.
BGE-M3 returns cosine scores; XSimGCL returns dot-product scores on
different scales.
MinMax normalisation maps each independently to $[0,1]$ before fusion;
the MLP outputs only scalar weights---not a mixed vector.
Linear aggregation is used instead of Softmax over final scores because
Softmax would exponentially exaggerate small numeric differences and
destroy inter-tower calibration.

\textbf{Missing Graph Edge Handling.}
When item $i$ has no learned XSimGCL embedding (added post-training),
its CF score is set to 0.
A warm user with $(\alpha,\beta)=(0.6,0.3)$ then obtains
$s(i)=0.6\,s_\text{sem}(i)$---automatically falling back to
semantic-only without any conditional branch in the code.

The fused score is:
$s(i) = \alpha\,\hat{s}_\text{sem}(i) + \beta\,\hat{s}_\text{CF}(i)
       + \gamma\,\hat{s}_\text{cooc}(i)$,
where $\hat{s}(\cdot)=\mathrm{minmax}(s(\cdot))$.

\subsection{Item-Item Co-Occurrence}

For every pair $(a,b)$ co-liked by at least $m_\text{min}=3$ distinct
users:
\begin{equation}
  \mathrm{cooc}(a,b)=\frac{|\mathcal{U}_a\cap\mathcal{U}_b|}
                          {\sqrt{|\mathcal{U}_a|\cdot|\mathcal{U}_b|}}.
\end{equation}
Top-50 neighbours per item are persisted in MongoDB and loaded in-memory
at startup.
This signal grows autonomously with the user base, providing a
self-improving audience-taste channel that IMDb/TMDB cannot supply.

\subsection{Bayesian Quality Filtering}

Language-specific IMDb weighted rating with prior vote threshold $m_\ell$:
\begin{equation}
  \mathrm{WR}(i) = \frac{v_i}{v_i+m_\ell}R_i
               + \frac{m_\ell}{v_i+m_\ell}C,
  \label{eq:bayes}
\end{equation}
$C=7.0$ (global mean), $m_\ell$ per Table~\ref{tab:votethresh}.
Scores are normalised \emph{within each language group} so regional films
compete fairly against English titles.
A stricter credibility penalty ($m=20{,}000$ regional; $m=50{,}000$
global) further adjusts inflated low-vote ratings
(e.g.\ a film rated 9.1 with 3,000 votes adjusts to $\approx$7.1).

\begin{table}[h]
\caption{Language-Specific Vote Thresholds for Eq.~(\ref{eq:bayes})}
\label{tab:votethresh}
\centering\small
\begin{tabular}{ll|ll}
\toprule
\textbf{Language} & $m_\ell$ & \textbf{Language} & $m_\ell$ \\
\midrule
Telugu    & 300  & Japanese  & 1000 \\
Tamil     & 500  & Korean    & 1000 \\
Malayalam & 500  & European  & 2000 \\
Hindi     & 800  & English   & 3000 \\
\bottomrule
\end{tabular}
\end{table}

\subsection{Genre-Dominance Reranking}

Genre dominance:
$D(i,G_u) = |G_i\cap G_u|/(|G_i|+1)$.
Total genre-fit score:
\begin{align}
  f_\text{genre}(i) &= 0.35\,D(i,G_\text{sel}) \nonumber\\
  &\quad + 0.25\,\sigma(n_\text{liked})\cdot\mathbf{1}[G_i\cap G_\text{liked}] \nonumber\\
  &\quad + 0.08\cdot\mathbf{1}[G_i\cap G_\text{okay}] \nonumber\\
  &\quad - 0.20\cdot\mathbf{1}[G_i\cap G_\text{disliked}] \nonumber\\
  &\quad - 0.10\cdot\mathbf{1}[G_i\cap G_\text{skipped}],
\end{align}
$\sigma(n)=1/(1+0.3n)$ prevents a frequently-liked genre from dominating.
Results interleave pure-genre ($D\geq0.70$) and mixed-genre items,
and relax filtering when the pool drops below 30 candidates.

\subsection{Two-Stage DPP Diversity Reranking}

\textbf{Stage 1 --- Categorical DPP.}
Feature matrix $F\in\mathbb{R}^{n\times(|G|+|L|)}$ of genre and language
one-hot vectors (weights $w_g=0.6$, $w_l=0.4$; reduced to $w_g=0.1$
when user has explicitly selected genres).
Greedy MAP selects $K=90$ from top-200.

\textbf{Stage 2 --- Embedding DPP.}
Quality-similarity kernel:
$L_{ij}=q_i q_j\,\mathbf{e}_i^\top\mathbf{e}_j$,
$q_i=\sqrt{s(i)}\cdot\delta_\text{lang}(i)$,
where $\delta_\text{lang}=1.3$ for non-dominant-language items.
Greedy MAP iterates $i^*=\arg\max_i\det(L_{S\cup\{i\}})/\det(L_S)$,
selecting $K=50$ from top-150.

% ─────────────────────────────────────────────────────────────
\section{System Design}

\textbf{Frontend.}
A Next.js application on Vercel (\url{onlymovies.app}) handles user
authentication, a 24-movie onboarding slate (8 per page; 70\%
MovieLens-linked, IMDb$\,\geq6.0$ / votes$\,\geq3{,}000$),
swipe-based interaction logging, and the personalised recommendation feed.

\textbf{Backend.}
The system was initially validated via a Gradio prototype to confirm
model inference correctness before engineering the production backend.
The final backend is a \textbf{Dockerised FastAPI container} deployed on
Hugging Face Spaces.
At startup the container downloads FAISS indexes and XSimGCL embeddings
from the \texttt{ml8r/cinematch} Hugging Face Dataset.
FAISS \texttt{IndexFlatIP} was preferred over IVF/HNSW variants for
correctness: exact inner-product search on $L_2$-normalised vectors
equals cosine similarity with zero approximation error.
The IDMap2 wrapper stores TMDB IDs directly, removing any secondary
lookup table.

\textbf{Database.}
MongoDB Atlas stores: user documents (profile, \texttt{cf\_embedding},
\texttt{is\_warm}), session state, and interaction logs
(\texttt{user\_id}, \texttt{tmdb\_id}, \texttt{action},
\texttt{context}, \texttt{timestamp}).
The \texttt{item\_cooccurrence} collection is rebuilt nightly and
loaded in-memory at container start.

% ─────────────────────────────────────────────────────────────
\section{Experimental Design}
\label{sec:expdesign}

\textbf{Hardware.}
XSimGCL training: UF HiPerGator NVIDIA B200.
Embedding generation and baseline runs: NVIDIA A100 40\,GB.

\textbf{Framework.}
RecBole-GNN~\cite{recbolegnn} (PyTorch 2.x) for all CF models.
FAISS 1.7 for indexing.

\textbf{Dataset Splits.}
ML-32M split temporally (80/10/10 per user).
Stratified held-out set: 979 test users across 5 activity bands
(5--50 / 50--100 / 100--200 / 200--500 / 500+;
183/196/200/200/200 users; 10 held-out items per user at $\geq4.0$\,stars;
9,790 test pairs total).

\textbf{Protocol.}
RecBole \texttt{uni100}: each test item ranked against 100 sampled negatives.
Metrics: NDCG@$K$, Recall@$K$, MRR@$K$, $K\in\{10,20,50\}$.

\textbf{ML-1M Validation.}
XSimGCL validated on ML-1M (6,040 users; 3,629 items; 836,478
interactions; 96.18\% sparse) under full-rank evaluation to confirm
implementation correctness before committing full ML-32M GPU hours.

% ─────────────────────────────────────────────────────────────
\section{Results and Analysis}
\label{sec:results}

\subsection{Main Comparison on ML-32M}

\begin{table}[h]
\caption{Test Performance --- MovieLens-32M (uni100, $n=979$ users)}
\label{tab:main}
\centering\small
\begin{tabular}{lccccc}
\toprule
\textbf{Method} & \textbf{N@10} & \textbf{R@10} & \textbf{M@10}
                & \textbf{N@20} & \textbf{R@20}\\
\midrule
LightGCN  & 0.6410 & 0.6124 & 0.7312 & 0.6687 & 0.7201\\
NCL       & 0.6890 & 0.6754 & 0.7798 & 0.7143 & 0.7851\\
SimGCL    & 0.6920 & 0.6802 & 0.7834 & 0.7181 & 0.7913\\
\midrule
\textbf{XSimGCL} & \textbf{0.7299} & \textbf{0.7277} & \textbf{0.8045}
                 & \textbf{0.7564} & \textbf{0.8298}\\
\midrule
\multicolumn{6}{l}{\footnotesize $\Delta$ vs LightGCN:
  $+13.9\%$ / $+18.8\%$ / $+10.0\%$ (N/R/M@10)}\\
\multicolumn{6}{l}{\footnotesize $\Delta$ vs SimGCL:\quad
  $+5.5\%$ / $+7.0\%$ / $+2.7\%$}\\
\bottomrule
\multicolumn{6}{l}{\footnotesize N=NDCG, R=Recall, M=MRR}
\end{tabular}
\end{table}

XSimGCL's noise-contrastive augmentation consistently lifted all
metrics above all baselines.
The MRR@10 gain of $+10.0\%$ over LightGCN (0.8045 vs.\ 0.7312) is
especially important for recommendation UIs where only the first few
positions receive user attention.

\subsection{CF Backbone Comparison on ML-1M}

Table~\ref{tab:ml1m} confirms implementation correctness against nine
baselines under full-rank evaluation (all items, no negative sampling).
Full-rank scores are lower than uni100 by design and are not directly
comparable to Table~\ref{tab:main}.

\begin{table}[h]
\caption{CF Backbone Comparison on ML-1M --- Full-Rank Evaluation}
\label{tab:ml1m}
\centering\scriptsize
\begin{tabular}{lccccc}
\toprule
\textbf{Method} & \textbf{N@10} & \textbf{R@10} & \textbf{M@10}
                & \textbf{Hit@10} & \textbf{P@10}\\
\midrule
BPR      & 0.2401 & 0.1776 & 0.4187 & 0.7199 & 0.1779\\
NeuMF    & 0.2271 & 0.1651 & 0.4020 & 0.7029 & 0.1700\\
NGCF     & 0.2508 & 0.1814 & 0.4354 & 0.7239 & 0.1850\\
LightGCN & 0.2538 & 0.1861 & 0.4388 & 0.7330 & 0.1863\\
LightGCL & 0.2479 & 0.1867 & 0.4283 & 0.7370 & 0.1815\\
SGL      & 0.2505 & 0.1889 & 0.4315 & 0.7392 & 0.1843\\
HMLET    & 0.2490 & 0.1847 & 0.4297 & 0.7305 & 0.1836\\
NCL      & 0.2702 & 0.2021 & 0.4599 & 0.7565 & 0.1962\\
SimGCL   & 0.2667 & 0.2029 & 0.4550 & 0.7640 & 0.1933\\
\midrule
\textbf{XSimGCL} & \textbf{0.2750} & \textbf{0.2116}
  & \textbf{0.4638} & \textbf{0.7743} & \textbf{0.1987}\\
\bottomrule
\end{tabular}
\end{table}

\subsection{Training Convergence}

NDCG@20 improved from 0.7083 (epoch~2) to 0.7901 (epoch~47), with most
gain in the first 20 epochs (0.7791).
BPR loss: $40.66\to3.34$; contrastive loss: $0.001\to1.18$;
regularisation: $89.27\to81.27$.

\subsection{Statistical Significance}

\begin{table}[h]
\caption{Statistical Significance --- XSimGCL vs.\ Baselines
(Paired $t$-test, one-sided, $n=979$, $\sigma_\Delta\approx0.22$)}
\label{tab:stats}
\centering\small
\begin{tabular}{lccl}
\toprule
\textbf{Comparison} & $t$ & $p$ & \textbf{95\% CI}\\
\midrule
vs.\ LightGCN & 12.65 & $<0.001$ & $[+0.075,+0.103]$\\
vs.\ NCL      & 5.82  & $<0.001$ & $[+0.027,+0.055]$\\
vs.\ SimGCL   & 5.39  & $<0.001$ & $[+0.024,+0.052]$\\
\bottomrule
\end{tabular}
\end{table}

\subsection{CF Backbone Ablation}

\begin{table}[h]
\caption{CF Backbone Ablation (NDCG@10, ML-32M, uni100)}
\label{tab:ablation}
\centering\small
\begin{tabular}{lcc}
\toprule
\textbf{Configuration} & \textbf{NDCG@10} & \textbf{R@10}\\
\midrule
ALS (Week 4, initial)    & 0.4821 & 0.4703\\
LightGCN                 & 0.6410 & 0.6124\\
UltraGCN (Week 8)        & 0.6987 & 0.6871\\
SimGCL                   & 0.6920 & 0.6802\\
BGE-M3 only (FAISS)      & 0.6341 & 0.6053\\
\textbf{XSimGCL (final)} & \textbf{0.7299} & \textbf{0.7277}\\
\bottomrule
\end{tabular}
\end{table}

Removing the CF backbone entirely (BGE-M3 only) drops NDCG@10 by 9.0\%
relative (0.7299$\to$0.6341), quantifying the value of collaborative
signals.
ALS's 51\% gap validates the GNN backbone choice.

\subsection{Cross-Cultural Discovery Rate}

\begin{table}[h]
\caption{CCDR --- Live User Study ( English-history users)}
\label{tab:ccdr}
\centering\small
\begin{tabular}{lcc}
\toprule
\textbf{System} & \textbf{CCDR} & \textbf{95\% CI}\\
\midrule
XSimGCL CF only (no DPP) &  9\% & $[6\%,12\%]$\\
CineMatch Full (+ DPP)   & 34\% & $[26\%,42\%]$\\
\bottomrule
\end{tabular}
\end{table}

The 3.8$\times$ increase ($9\%\to34\%$) confirms the combined effect of
BGE-M3's cross-lingual embedding alignment and the DPP's $1.3\times$
language-diversity bonus.

% ─────────────────────────────────────────────────────────────
\section{Error Analysis \& Interpretability}
\label{sec:erroranalysis}

\textbf{Score Decomposition (XAI).}
Every recommendation exposes three interpretable scores:
$\hat{s}_\text{sem}$ (BGE-M3 content similarity),
$\hat{s}_\text{CF}$ (XSimGCL behavioural match), and
$s$ (final fused), alongside language, genre, and quality tags.
This \emph{transparency-by-design} fulfils Microsoft HAI Guideline~G11
(``Make clear why the system did what it did''):
a user can inspect that a Korean thriller appeared because content
similarity was 0.87 while CF was 0.21, indicating the system acted on
narrative resemblance rather than population behaviour.

\textbf{HAI Trust and Safety.}
When the system's candidate pool falls below 30 items (indicating
unusual taste profiles), genre and vote filters are relaxed and the
system surfaces a diversity notice (G10: ``Scope services when in doubt'').
Cold-start users are explicitly shown onboarding context:
``Based on your 8 rated films, here are suggestions---rate more for
personalised results.''

\textbf{Failure 1 --- Genre Collapse (resolved).}
After $\approx$40 interactions, users reported ``I only get action films.''
\emph{Root cause}: hard genre filter plus flat 3,000-vote floor exhausted
the non-action international film pool.
\emph{Fix}: soft dominance scoring; language-specific vote floors; pool
floor (min 30 candidates) dynamically relaxes both filters.

\textbf{Failure 2 --- Cold-Start Irrelevance (resolved).}
Users with $<5$ rated movies received inconsistent recommendations
because single-movie BGE-M3 queries were insufficient context.
\emph{Fix}: onboarding collects 8--24 ratings before personalised output.

\textbf{Failure 3 --- Language Embedding Mismatch (partially resolved).}
Korean films described entirely in English in TMDB embedded closer to
English-language thrillers than to Korean-language cinema.
\emph{Fix}: explicit language-aware reranking ($3\times$ boost target
language; $0.3\times$ suppression for English items).

\textbf{Limitations.}
MLP gating is hand-initialised; gradient training is feasible as logs
accumulate.
XSimGCL is transductive: items added post-training receive no CF score.
The BGE-M3 inductive path mitigates this; a Feature-Initialised GNN
fold-in mechanism would close the gap permanently.

% ─────────────────────────────────────────────────────────────
\section{Challenges and Lessons}
\label{sec:challenges}

\textbf{Computational Scale.}
Embedding 1.37M movies with BGE-M3 and training on 20M interactions
both exceeded local machine memory.
Checkpointing every 1,000 batches prevented costly restarts after
HiPerGator pre-emptions.
Validating on ML-1M first (each epoch $\approx$8\,s vs.\ $\approx$107\,s
on ML-32M) saved $\approx$30\,GPU-hours.

\textbf{TMDB Vote Sparsity.}
Many acclaimed international films had zero TMDB votes.
Joining IMDb crowd-sourced counts on \texttt{imdb\_id} resolved this
for $\approx$738K titles, requiring careful deduplication by
\texttt{imdb\_id} fuzzy match.

\textbf{Encoder Index Selection.}
The Qwen3-4B index (14\,GB) produced marginally higher NDCG on niche
queries but was impractical at deployment.
BGE-M3 (5.6\,GB) provided fast, accurate retrieval within budget.
This trade-off was only discoverable by building and testing both indexes.

\textbf{Offline-Online Gap.}
Optimising NDCG@10 offline did not prevent genre collapse or cold-start
degradation in live deployment.
Real users surfaced failures (``it stops showing interesting movies
after a while'') that held-out test sets could not detect, validating
live deployment as a critical complement to offline metrics.

% ─────────────────────────────────────────────────────────────
\section{Conclusion}
\label{sec:conclusion}

CineMatch demonstrated that cross-cultural movie discovery is achievable
without generative LLMs or synthetic data.
A five-stage pipeline---late fusion of XSimGCL (NDCG@10\,=\,0.7299) and
BGE-M3, MLP gating, Bayesian quality scoring, genre-dominance reranking,
and two-stage DPP diversity---outperformed LightGCN by 13.9\% and
SimGCL by 5.5\% on NDCG@10 (both $p<0.001$), and increased
cross-cultural exposure 3.8$\times$ in live testing.
The language-specific Bayesian vote prior and the two-stage DPP
are the primary novelties over existing graph-CF pipelines; together
they elevate regional cinema from invisible to discoverable.
Real-user feedback drove two architecture corrections that no offline
metric would have surfaced.

\textbf{Future work:}
(1)~\emph{Feature-Initialised Inductive GNN}: replace transductive
\texttt{nn.Embedding} ID lookup with a fixed BGE-M3 feature matrix;
train only graph aggregation weights so new movies are foldable at
inference time without full graph retraining (user nodes initialised
by mean-pooling their interacted item embeddings).
(2)~End-to-end MLP gate training via REINFORCE once interaction logs
are sufficient.
(3)~RAG conversational layer: map natural-language queries through an
LLM into BGE-M3 query vectors for dialogue-driven discovery.
(4)~Empirical selection between the three built FAISS indexes as
hosting infrastructure scales.

% ─────────────────────────────────────────────────────────────
\begin{thebibliography}{9}

\bibitem{xsimgcl}
J.~Yu, H.~Yin, X.~Xia, T.~Chen, L.~Cui, and Q.~V.~H.~Nguyen,
``Are Graph Augmentations Necessary? Simple Graph Contrastive
Learning for Recommendation,''
\textit{Proc.\ ACM SIGIR}, 2022. arXiv:2209.02544.

\bibitem{bgem3}
J.~Chen, S.~Xiao, P.~Zhang, K.~Luo, D.~Lian, and Z.~Liu,
``BGE M3-Embedding: Multi-Lingual, Multi-Functionality,
Multi-Granularity Text Embeddings,''
arXiv:2402.03216, 2024.

\bibitem{lightgcn}
X.~He, K.~Deng, X.~Wang, Y.~Li, Y.~Zhang, and M.~Wang,
``LightGCN: Simplifying and Powering Graph Convolution
Network for Recommendation,''
\textit{Proc.\ ACM SIGIR}, pp.\ 639--648, 2020.

\bibitem{ncl}
Z.~Lin, C.~Tian, Y.~Hou, and W.~X.~Zhao,
``Improving Graph Collaborative Filtering with
Neighbourhood-Enriched Contrastive Learning,''
\textit{Proc.\ WWW}, pp.\ 2320--2329, 2022.

\bibitem{movielens}
F.~M.~Harper and J.~A.~Konstan,
``The MovieLens Datasets: History and Context,''
\textit{ACM Trans.\ Interact.\ Intell.\ Syst.}, vol.~5, no.~4, 2015.

\bibitem{dpp}
L.~Chen, G.~Zhang, and E.~Zhou,
``Fast Greedy MAP Inference for Determinantal Point Process
to Improve Recommendation Diversity,''
\textit{Proc.\ NeurIPS}, 2018.

\bibitem{recbolegnn}
Y.~Wang et al.,
``Towards Unified Graph Neural Networks for Recommendation,''
arXiv:2206.07089, 2022.

\bibitem{qwen3}
Qwen Team, Alibaba Cloud,
``Qwen3 Technical Report,''
arXiv:2505.09388, 2025.

\bibitem{sasrec}
W.~Kang and J.~McAuley,
``Self-Attentive Sequential Recommendation,''
\textit{Proc.\ IEEE ICDM}, pp.\ 197--206, 2018.

\bibitem{bert4rec}
F.~Sun, J.~Liu, J.~Wu, C.~Pei, X.~Lin, W.~Ou, and H.~Jiang,
``BERT4Rec: Sequential Recommendation with
Bidirectional Encoder Representations from Transformers,''
\textit{Proc.\ ACM CIKM}, pp.\ 1441--1450, 2019.

\bibitem{rag}
P.~Lewis et al.,
``Retrieval-Augmented Generation for
Knowledge-Intensive NLP Tasks,''
\textit{Proc.\ NeurIPS}, 2020.

\bibitem{deepfm}
H.~Guo, R.~Tang, Y.~Ye, Z.~Li, and X.~He,
``DeepFM: A Factorization-Machine based Neural Network
for CTR Prediction,''
\textit{Proc.\ IJCAI}, pp.\ 1725--1731, 2017.

\end{thebibliography}

% ─────────────────────────────────────────────────────────────
\appendix
\section*{Appendices}

\textbf{A. GitHub Repository.}\\
\url{https://github.com/roop-yaswanth/CineMatch}\\
\texttt{README.md}, \texttt{requirements.txt}, \texttt{app.py}
($\approx$5,800 lines), \texttt{cinematch-web/} (Next.js),
\texttt{RecBole-GNN/} (training), \texttt{models/} (FAISS manifests,
XSimGCL config).

\vspace{1mm}
\textbf{B. Live Demo.}\\
\url{https://onlymovies.app}
(Next.js, Vercel; Docker FastAPI backend: Hugging Face Spaces).\\
2-minute video walkthrough:
\url{https://github.com/roop-yaswanth/CineMatch}

\vspace{1mm}
\textbf{C. Compute Resources.}\\
Training: NVIDIA B200, UF HiPerGator (\texttt{egn6933/nagabhairava.r}).\\
Embedding \& baselines: NVIDIA A100 40\,GB.\\
Total GPU-hours (training + experimentation + indexing):
$\approx$14\,h.

\vspace{1mm}
\textbf{D. Reproducibility.}\\
Seed: 42 (RecBole \texttt{reproducibility: true};
CUDA deterministic enabled).
FAISS index construction is deterministic given fixed input order.

\end{document}
